\section{An exact comparison near a common diameter}\label{sec:local}
\begin{theorem}[Explicit independent local comparison]\label{thm:local}
Set
\begin{equation}\label{eq:local-constants}
 \mathcal A=e^{10000},\qquad\lambda=e^{-\mathcal A},
                   \qquad N=\lceil e^{\mathcal A}\rceil.
\end{equation}
Every independent array of the form
\eqref{eq:clusters}--\eqref{eq:extra-bounds}, with $m\ge N$, has
$\mathcal R\le\be$. Its ratio uses the total variance and the sum
of third absolute moments of all the original coordinates.
\end{theorem}

We follow the two residual regimes of
\citet[Section~3]{HeCheng2026}. The estimates below are uniform in
unequal probabilities, unequal gaps, and the additional coordinates.

\subsection{Normalization and conditional estimates}
Put $v_i=p_i(1-p_i)$ and define the variance-weighted mean gap
\[
 \bar h=\frac{\sum_i v_i h_i}{\sum_i v_i}.
\]
Then $(1-\lambda)h\le\bar h\le(1+\lambda)h$.
Divide every coordinate by $\bar h$, retaining $h_i,U_i,Z_j$ for the
rescaled gaps and residuals. Write
\begin{gather*}
 V=\sum_i v_i,\quad\mu=\sum_i p_i,
 \quad h_i=1+a_i,\quad\sum_i v_i a_i=0,\\
 K_0=\sum_iB_i,\quad A=\sum_i a_i(B_i-p_i),\quad
 U=\sum_iU_i+\sum_j Z_j,\quad W=A+U,\quad S=K_0+W,\\
 \Lambda_A=\sum_i v_i a_i^2,\quad
 \Lambda_U=\sum_i\E U_i^2+\sum_j\E Z_j^2,\quad
 \Lambda=\Lambda_A+\Lambda_U,\quad B^2=V+\Lambda.
\end{gather*}
The uncentered sum $S$ has mean $\mu$, and $\Cov(K_0,W)=0$.
Put $\eps_0=\max(\max_i\|U_i\|_\infty,\max_j\|Z_j\|_\infty)$,
with an empty maximum zero, and $\xi=\Lambda/V$.
The hypotheses imply
\begin{equation}\label{eq:local-domains}
 3m/16\le V\le m/4,\quad\max_i|a_i|\le3\lambda,\quad
 \eps_0\le2\lambda,\quad\xi\le8\lambda,\quad\Lambda\le2\lambda m.
\end{equation}
For example the upper bound before enlargement is
$\xi\le[\lambda+(7/4)\lambda^2]/[(3/16)(1-\lambda)^2]$.

Let $T_X$ and $\kappa$ be the total absolute and signed third central
moments of the original array. Its reference Bernoulli moments are
\[
 T_0=\sum_i v_i(p_i^2+(1-p_i)^2),\qquad
 \kappa_0=\sum_i v_i(1-2p_i).
\]
Conditional centering gives
\begin{align}
 |T_X-T_0|,\ |\kappa-\kappa_0|
      &\le2\sqrt{m\Lambda_A}+4\Lambda,                  \label{eq:moment-budget}\\
 T_X\ge V/3,\quad1/4\le T_X/B^2&\le2,\quad
 B^2\le m,\quad B/2\le P:=B^3/T_X\le4B.                \label{eq:prefactor}
\end{align}
Indeed $|(1+a)^3-1|\le4|a|$ and
$\sum v_i|a_i|\le\sqrt{V\Lambda_A}$.
Within each cluster the sign is fixed: its unperturbed magnitude is
at least $7/32>\eps_0$. Expanding the cubic changes either third
moment by at most $4\E U_i^2$, since the linear term vanishes.
An extra coordinate contributes at most $\eps_0\E Z_j^2$.
Jensen gives $T_X\ge(7/8)^3T_0\ge343V/1024>V/3$.
Every original coordinate has absolute value at most two, so
$T_X\le2B^2$, proving \eqref{eq:prefactor}.

For $m\ge10^{12}$ and $|k-\mu|\le6\sqrt m$,
Appendix~\ref{app:conditional} proves
\begin{align}
 |m_k|:=|\E(A\mid K_0=k)|&\le10^4\sqrt{\Lambda_A}/\sqrt m,
                                                        \label{eq:conditional-mean}\\
 |\Var(A\mid K_0=k)-\Lambda_A|&\le40000\Lambda_A/\sqrt m,\notag\\
 \E[(A-m_k)^4\mid K_0=k]&\le10^7\Lambda_A^2.             \label{eq:conditional-gap}
\end{align}
It also proves the relative label bound
\begin{equation}\label{eq:conditional-label}
 \left|\frac{\Pp(B_i=b\mid K_0=k)}{\Pp(B_i=b)}-1\right|
                         \le500/\sqrt m\qquad(b=0,1).
\end{equation}
For $\Lambda>0$ define
\begin{equation}\label{eq:DH}
 D=\Lambda^2+\eps_0^2\Lambda_U,\qquad
 H=\max\left\{\sqrt{\Lambda_A},
                   \frac\Lambda{\sqrt{\Lambda+\eps_0^2}}\right\}.
\end{equation}
Then, in the same count range,
\begin{gather}
 \E(W\mid K_0=k)=m_k,\quad
 \Lambda/2\le\Var(W\mid K_0=k)\le2\Lambda,\notag\\
 \E[(W-m_k)^4\mid K_0=k]\le10^8D,\quad
 \E(|W-m_k|\mid K_0=k)\ge H/30000,                       \label{eq:full-conditional}\\
 \E(W^4\mid K_0=k)\le10^9D,\qquad\E W^4\le400D.\notag
\end{gather}
To see the extension, condition first on all labels. The noise then
has variance $T_{\mathbf b}\le4\Lambda_U$ and fourth moment at most
$3T_{\mathbf b}^2+\eps_0^2T_{\mathbf b}$.
Its mean given $K_0$ is zero, and \eqref{eq:conditional-label}
compares its variance to $\Lambda_U$. Combine these facts with
\eqref{eq:conditional-gap} and $|x+y|^4\le8(|x|^4+|y|^4)$.
Interpolation $\E|Y|\ge(\E Y^2)^{3/2}/(\E Y^4)^{1/2}$ supplies
the second term in $H$. Conditional Jensen given the labels supplies
the first from $A-m_k$. In particular $|m_k|\le10^4H/\sqrt m$
uniformly, even when $\Lambda_A$ is much smaller than $\eps_0^2/m$.
Finally $\max a_i^2\le16\Lambda_A/3$ gives
$\E A^4\le9\Lambda_A^2$, and hence the unconditional bound.

\subsection{Uniform approximation and central thresholds}
Let $J$ be the CDF of $S$ plus an independent uniform variable on
$[-1/2,1/2]$. If $T\ge10$ and $\xi(T+1)^2\le1/56$, put
\begin{align}
 E_m(T)={}&\frac{401(1+\log T)}{\sqrt m}
       +(1+\log T)\xi(T+1)^2\notag\\
       &+60000(1+\log T)\sqrt m\,e^{-m/262144}+56/T.       \label{eq:local-error}
\end{align}
Appendix~\ref{app:fourier} proves, for $z=(x-\mu)/B$,
\begin{equation}\label{eq:local-jitter}
 \sqrt m\left\|J(x)-\Phi(z)
             -\frac\kappa{6B^3}(1-z^2)\phi(z)\right\|_\infty\le E_m(T).
\end{equation}
The expansion includes the third moments of the additional coordinates.

Fix
\[
 a_0=e^{1000},\quad\mathcal C=e^{1004},\quad\mathcal B=e^{1006},
 \quad q=e^{-\mathcal C},\quad T=e^{\mathcal B}.
\]
The inequalities $\mathcal A>100\mathcal B$,
$\mathcal B>7\mathcal C$, and $\mathcal C>50a_0$ imply, for every
$m\ge N$,
\begin{equation}\label{eq:local-numerical}
 \xi(T+1)^2<1/56,\quad E_m(T)<q/1000,\quad
 m^{-1/2}<q/1000,\quad200\lambda<q/1000.
\end{equation}
For a direct check, the four terms in \eqref{eq:local-error} are at most
\[
 802\mathcal B e^{-\mathcal A/2},\quad
 64\mathcal B e^{-\mathcal A+2\mathcal B},\quad
 240000\mathcal B(262144)^2e^{-3\mathcal A/2},\quad56e^{-\mathcal B}.
\]
Each is below $q/4000$. The third uses $e^{-x}\le2/x^2$, so this
verification places no upper bound on $m$.

Shift the jitter by $\pm1/2$. Either normalized signed discrepancy is
at most
\begin{equation}\label{eq:local-envelope}
 \phi(z)\left\{\frac{B^2}{2T_X}
      +\frac{|\kappa|}{6T_X}|1-z^2|\right\}+4E_m(T)+m^{-1/2}.
\end{equation}
The normalized Taylor error is at most $1/(4B)<m^{-1/2}$, using
$\|\phi'\|_\infty<1/4$,
$\|((1-z^2)\phi(z))'\|_\infty<5/4$, and \eqref{eq:prefactor}.
The main term is at most $(3+z^2)\phi(z)$, below $1/100$ for
$|z|\ge4$. Thus the full bound there is below $1/24$.

\subsection{Small residual variance}\label{subsec:small-residual}
Suppose $0<\Lambda\le e^{-1000}$. If $\Lambda=0$, the result is
exactly \eqref{eq:finite-input}. We use a one-sided loss estimate:
if $\E Y=0$, $\E Y^4\le1/160000$, $f(0)=0$, and
$3/4\le f'\le5/4$ on $[-7/10,7/10]$, then for $|u|\le3/5$,
\begin{align}
 \Pp(Y>u)+f(u)&\ge\tfrac38\E|Y|-176000\E Y^4,\notag\\
 \Pp(Y<u)-f(u)&\ge\tfrac38\E|Y|-176000\E Y^4.             \label{eq:one-sided}
\end{align}
This is the clipping argument of
\citet[Lemma~3.2 and equation~(83)]{HeCheng2026}. Appendix~\ref{app:conditional}
gives its proof with these constants and the stated threshold range.

The count estimate in Appendix~\ref{app:fourier} gives
\begin{equation}\label{eq:central-count}
 b_k:=\Pp(K_0=k)\ge e^{-100}/\sqrt m
                         \qquad(|k-\mu|\le6\sqrt m).
\end{equation}
Write a threshold with $|z|<4$ as $k+u$, $|u|\le1/2$, where $k$
is a nearest integer. Then $|k-\mu|\le5\sqrt m$.
For $f_k(u)=[\Phi((k-\mu+u)/B)-\Phi((k-\mu)/B)]/b_k$,
we have $3/4\le f_k'\le5/4$ on $[-1,1]$.
Indeed the relative count error is at most $18e^{100}/\sqrt m$,
and the logarithm of the Gaussian density ratio is at most
$100\xi+30/\sqrt m$; both are below $1/100$.

The exact count decompositions are
\begin{align}
 F_S(k+u)&=F_{K_0}(k)-b_k\Pp(W>u\mid K_0=k)+e^+_{k,u},\notag\\
 F_S^-(k+u)&=F_{K_0}(k-1)+b_k\Pp(W<u\mid K_0=k)+e^-_{k,u}.
                                                        \label{eq:layers}
\end{align}
Each remainder is bounded in absolute value by
\[
 \sum_{j\ne k}b_j\Pp\{|W|\ge|j-k|-1/2\mid K_0=j\}
                         \le2\cdot10^{11}D/\sqrt m.
\]
For $|j-k|\le\sqrt m$, use $b_j\le3/\sqrt m$,
\eqref{eq:full-conditional}, and
$\sum_{r\ne0}(|r|-1/2)^{-4}<40$.
For more distant counts use $\E W^4\le400D$ and a denominator
at least $m^2/16$.

Apply \eqref{eq:one-sided} to $Y=W-m_k$ conditionally on $K_0=k$,
using $f_k(v+m_k)-f_k(m_k)$ and $v=u-m_k$.
The bounds above give $|m_k|<1/20$ and the required fourth-moment
bound. Also $e^{-100}/5\le Pb_k\le5$.
The resulting normalized loss is at least $e^{-120}H$.
Its fourth-moment and leakage errors are below $10^{14}D$, and its
conditional-mean error is at most $62500\sqrt{\Lambda_A}/\sqrt m$.

Let $R_0=\mathcal R(B_1-p_1,\ldots,B_m-p_m)\le\be$.
Equation~\eqref{eq:moment-budget} gives
\[
 \left|\frac{B^3/T_X}{V^{3/2}/T_0}-1\right|
                  \le32\sqrt{\Lambda_A}/\sqrt m+96\Lambda/m.
\]
The normal variance change costs at most $\xi/8$ before multiplication
by $P$. Combining these errors with \eqref{eq:layers}, either central
normalized branch is at most
\begin{equation}\label{eq:small-comparison}
 R_0-e^{-120}H+e^{40}
           \left\{D+\frac{\sqrt{\Lambda_A}+\Lambda}{\sqrt m}\right\}.
\end{equation}
The full definition of $H$ gives
\[
 D/H\le(\Lambda+\eps_0^2)^{3/2},\quad
 \sqrt{\Lambda_A}/H\le1,\quad
 \Lambda/H\le\sqrt{\Lambda+\eps_0^2}.
\]
Since $\eps_0\le e^{-1000}$ and $m\ge e^{1000}$, the error is below
$e^{40}(3e^{-1500}+2e^{-500})H<e^{-120}H/2$.
Thus the central comparison is strict for every positive residual
variance. Finally $R_0\ge\phz/4>1/12$: apply the count estimate at
an integer nearest $\mu$ and take half its jump. The bound $1/24$
at distant thresholds is also below $R_0$.

\subsection{Positive residual variance and interval mass}
\label{subsec:large-residual}
Suppose $\Lambda\ge e^{-1000}$ and $|z|\le4$.
Choose all integers $|k-x|\le\max(1,\sqrt\Lambda)$. There are at
least $\sqrt\Lambda$ of them, all within $6\sqrt m$ of $\mu$.
Tilt the Bernoulli probabilities exponentially to make their count
mean $k$, and project the gap residual off the count, as in
Appendix~\ref{app:conditional}. On $K_0=k$ the residual is
$m_k^0+R_k$, where under the tilted measure $R_k$ is a sum of
independent centered pair residuals and additional coordinates, with
\begin{equation}\label{eq:tilted-residual}
 L:=\Var_\theta R_k\in[\Lambda/2,2\Lambda],\quad
 \Var_\theta K_0\ge m/10,\quad |m_k^0|\le\sqrt\Lambda,
 \quad\Cov_\theta(K_0,R_k)=0.
\end{equation}
Each residual summand is bounded by $8\lambda$.
The within-cluster noise variances change relatively by at most
$120/\sqrt m$ under the tilt; the extra coordinates do not change.
Thus these assertions include their actual variances.

Apply Lemmas~\ref{lem:conditional-fourier} and~\ref{lem:minorant}
to both quarter-intervals $(1/4,1/2)$ and $(-1/2,-1/4)$, with
\[
 L_0=\tfrac12e^{-1000},\quad C_0=4+2e^{500},\quad A_g=4e^{500},
 \quad\omega=16e^{A_g^2/6},\quad T_g=4\omega.
\]
The shift $y=x-k-m_k^0$ satisfies $|y|\le C_0\sqrt L$, and
$A_g\ge C_0+(7/16)/\sqrt{L_0}$.
The interval minorant has Gaussian expectation at least
$\phi(A_g)/(4\omega\sqrt L)$.
With $\beta=7/64$, the three terms in
\eqref{eq:conditional-fourier-error} are at most
\[
 128e^{-\mathcal A/2},\qquad128e^{500-\mathcal A},\qquad
                         10^7e^{4a_0-\mathcal A}.
\]
Their sum is below $e^{-10a_0}<\phi(A_g)/32$, because
$\phi(A_g)>e^{-9a_0}$, $\omega<e^{3a_0}$, and $T_g<e^{4a_0}$.
Also $8\lambda T_g<\beta/162$.
For the third term use $L\le m$, $e^{-x}\le2/x^2$, and $m\ge N$;
this is uniform over all possible residual variances.

Each conditional quarter-interval probability is at least
$\phi(A_g)/(8\omega\sqrt L)$. Multiply by the count probabilities
\eqref{eq:central-count} and sum over the chosen integers.
Since $L\le2\Lambda$ and $\mathcal C>50a_0$, this gives
\begin{equation}\label{eq:quarter-mass}
 \Pp\{S-x\in(1/4,1/2)\},\quad
 \Pp\{S-x\in(-1/2,-1/4)\}\ge q/\sqrt m.
\end{equation}
The two actual one-sided jitter losses therefore satisfy
\[
 J(x+1/2)-F_S(x)\ge q/(2\sqrt m),\qquad
 F_S^-(x)-J(x-1/2)\ge q/(2\sqrt m).
\]
After multiplication by $P\ge B/2$, each loss exceeds $q/10$.
Esseen's Bernoulli inequality gives $3V+|\kappa_0|\le\ce T_0$.
Using \eqref{eq:moment-budget}, the main term of
\eqref{eq:local-envelope} is consequently at most
\[
 \be+20\sqrt{\Lambda_A/m}+42\Lambda/m\le\be+200\lambda.
\]
Either central branch is below
$\be-q/10+200\lambda+4E_m(T)+m^{-1/2}<\be-q/20$.
The distant thresholds were already bounded by $1/24$.
This completes the proof of Theorem~\ref{thm:local}.
